\section{Distribución $\chi^2$}

La distribuci\'on $\chi^2$ aparece al sumar cuadrados de normales est\'andar. Su
tratamiento detallado se realiza en la secci\'on \ref{sec:3.9}; aqu\'i presentamos
los resultados esenciales.

\begin{definicion}[Distribuci\'on $\chi^2$]
Sean $Z_1, Z_2, \ldots, Z_\nu$ variables aleatorias independientes con $Z_i \sim N(0,1)$.
La variable aleatoria
\begin{align}
 \label{eq:3.2.11}
 \chi^2_\nu = \sum_{i=1}^\nu Z_i^2
\end{align}
sigue una distribuci\'on \emph{chi-cuadrada} con $\nu$ grados de libertad.
\end{definicion}

{Propiedades}
\begin{itemize}
 \item Media: $E(\chi^2_\nu) = \nu$.
 \item Varianza: $\text{Var}(\chi^2_\nu) = 2\nu$.
 \item Como caso particular de la familia gamma (secci\'on 2.11),
 $\chi^2_\nu = \text{Gamma}\!\left(\frac{\nu}{2}, 2\right)$.
 \item Reproductividad: si $\chi^2_{\nu_1}$ y $\chi^2_{\nu_2}$ son independientes,
 $\chi^2_{\nu_1} + \chi^2_{\nu_2} \sim \chi^2_{\nu_1 + \nu_2}$.
\end{itemize}

\begin{observacion}
La distribuci\'on $\chi^2$ surge naturalmente al estimar la varianza de una poblaci\'on
normal. Si $X_1, \ldots, X_n$ son normales iid y $S^2$ es la varianza muestral,
entonces
\begin{align}
 \frac{(n-1)S^2}{\sigma^2} \sim \chi^2_{n-1}.
\end{align}
Este resultado es la base de los intervalos de confianza para $\sigma^2$ y de la
prueba chi-cuadrada de bondad de ajuste.
\end{observacion}

\begin{definicion}[Densidad de la chi-cuadrada]
 \label{eq:5.3.1}
Como caso particular de la distribuci\'on gamma con $\alpha=\nu/2$ y $\beta=2$, la
funci\'on de densidad de $\chi^2_\nu$ es
\begin{align}
 f(x) = \frac{1}{2^{\nu/2}\Gamma(\nu/2)}\,x^{\nu/2-1}e^{-x/2}, \qquad x > 0.
\end{align}
\end{definicion}

\begin{teorema}[Teorema de Fisher]
 \label{eq:5.3.2}
Sean $X_1, \ldots, X_n$ variables aleatorias iid con $X_i \sim N(\mu, \sigma^2)$, con
media muestral $\bar{X}$ y varianza muestral $S^2$. Entonces:
\begin{enumerate}
 \item $\bar{X}$ y $S^2$ son \emph{independientes};
 \item $\dfrac{(n-1)S^2}{\sigma^2} \sim \chi^2_{n-1}$.
\end{enumerate}
\end{teorema}

\begin{observacion}[Idea de la demostraci\'on]
La independencia de $\bar{X}$ y $S^2$ es un resultado sorprendente: aunque $S^2$ se
calcula a partir de las desviaciones $(X_i-\bar{X})$, que dependen de $\bar{X}$, el
vector de desviaciones resulta ser estad\'isticamente independiente de $\bar{X}$
cuando la poblaci\'on es normal. La demostraci\'on formal usa una transformaci\'on
ortogonal del vector $(X_1,\ldots,X_n)$ que separa la direcci\'on de $\bar{X}$
(un solo grado de libertad) de las $n-1$ direcciones ortogonales restantes, cada una
capturando informaci\'on independiente de $\bar X$; la suma de cuadrados en esas
$n-1$ direcciones es exactamente $(n-1)S^2$, lo que explica tanto la independencia
como los $n-1$ grados de libertad (este es un caso particular del \emph{teorema de
Cochran}).
\end{observacion}

\begin{ejemplo}
 \label{exmp:5.3.1}
El contenido de las latas de refresco de una embotelladora sigue
$N(\mu, \sigma^2=4)$ (en mililitros, con $\sigma=2$). Se toma una muestra de $n=10$
latas y se calcula $S^2 = 7.2$. ¿Cu\'al es el valor del estad\'istico
$(n-1)S^2/\sigma^2$, y qu\'e tan extremo es este valor comparado con la distribuci\'on
$\chi^2_9$ (cuyo valor esperado es $9$)?
\end{ejemplo}

\begin{solucion}
El estad\'istico es
\begin{align}
 \frac{(n-1)S^2}{\sigma^2} = \frac{9 \cdot 7.2}{4} = \frac{64.8}{4} = 16.2.
\end{align}
Como $16.2 > 9 = E(\chi^2_9)$, la varianza muestral observada es mayor de lo
esperado bajo $\sigma^2=4$. Usando tablas de $\chi^2_9$, el valor cr\'itico al
$95\%$ es $\chi^2_{9,0.95} \approx 16.92$; como $16.2 < 16.92$, el valor no es lo
suficientemente extremo para rechazar $\sigma^2=4$ a nivel $\alpha=0.05$ (aunque
est\'a cerca del umbral).
\end{solucion}


